% ==========================================================================
%  Section 14 -- Reproduction.  Owner: W10.
%  Frozen label: sec:lean-reproduction.
%
%  Sources read for this section: build_all.py (steps 1-7), audit/count_lean.py,
%  audit/make_ledger.py, audit/make_bib.py, audit/verify_paper.py (the four checks
%  N/E/L and their negative controls, the docstring "known limits"), main.tex
%  (the two documented build commands) and ext/axiom_transcript.txt lines 8-18
%  (the two recorded commands).
%
%  Traps observed here.  (a) No build-job count and no wall-clock time is stated:
%  both are environment-dependent and would rot.  (b) The reproducibility claim is
%  limited to declaration names and frozen digests; the correspondence of a paper
%  sentence with a Lean statement is not machine-checked and is said so.
% ==========================================================================
\section{Reproduction}\label{sec:lean-reproduction}

Two entry points are supported and they are not equivalent. The paper's own gate
is \texttt{python build\_all.py}, run from the paper directory. In order it
regenerates the structural counts and both appendices -- the declaration ledger
from the sources, the axiom transcript from the recorded engine output -- runs the
negative controls of the paper checks and then the checks themselves, compiles
with Tectonic \TectonicVersion{}, refuses a PDF older than any source, and stamps
the digest of every source and of the output into \lmod{build/LAST_BUILD.txt}.
The four checks are: that every Lean name quoted in the body is a name the
generated ledger lists, or else a declaration of the package found in the Lean
sources; that the large structural counts in the body are macros rather than
typed values; that labels, cross-references, citations and bibliography
keys close and every section carries a label; and that the axiom story is
self-consistent, in the sense of \cref{sec:lean-contracts}. Stated precisely: the
count check rejects a typed literal only when the macro it should have used has
at least three digits or a value of at least forty, so a hand-typed
\texttt{two} would pass it, and the discipline for small numbers is editorial
rather than mechanical; the check also fails when a macro that the build computes
is quoted nowhere, which is the failure mode a typed count leaves behind. The
reference check covers every reference command the document uses, \verb|\eqref|
included. The compile step is followed by a scan of the log: a build that reports
success while leaving an overfull box or an undefined reference is a failed build
here, because the engine treats both as warnings and the reader sees both as
defects. The project around the paper carries \NumProofDocs{} proof documents,
\NumProbes{} verification probes and \NumFailureCards{} failure cards; none is
shipped with this paper and none is an input to any claim above. The second entry
point is the bare compile, \texttt{tectonic -C main.tex}, which builds the PDF
from whatever the generated files currently contain. It is what a reader who
wants only the document runs, and it proves nothing about the sources: a stale
generated file compiles as happily as a fresh one.

The footprints are reproduced in the separate project, in two steps, using the
recorded commands. First the declaring module is built, which populates the
object files that the second step reads. Then a short file is written whose
imports are the modules under audit and whose body is one
\texttt{\#print axioms <name>} line per declaration, and that file is elaborated
with \texttt{lake env lean}. The toolchain is Lean \LeanToolchain{}, the same one
the transcript header records. This is exactly the command whose output
\lmod{ext/axiom_transcript.txt} transcribes. The transcript was regenerated on
2026-11-03 by \lmod{audit/make_transcript.py}, on the sources in this repository,
and it is the artefact to compare against.  The digests of the statement contracts are in a different position and must not
be confused with the transcripts: they are recorded values in the per-node files
held beside the registry entries described in \cref{sec:lean-contracts}, and
because both contracts are marked remote-only the runner reports \lmod{SKIP}
before it computes a digest at all.  That flag governs the digest runner, which
points at the project holding the recorded values; it does not say that the
package cannot be elaborated here, and \cref{sec:lean-overview} and the paragraph
below record that it can.  Nothing shipped with this paper exercises the digest,
but a reader can run the build.

A reproducer should expect the following, and should treat any deviation as a
finding. The package builds on the revision it pins --- Lean \LeanToolchain{}
with Mathlib \texttt{85e3a25e}, the commit \lmod{lake-manifest.json} pins for the
tag \lmod{lakefile.lean} names --- and the toolchain it reports is the one
the transcript header records: the footprints are valid on that toolchain, and
comparing them across toolchains is not a comparison. Each
\texttt{\#print axioms} line prints the permitted set --- all
\NumChainTheorems{} of them do, in Appendix~B; a line carrying a name outside it,
or an empty footprint in place of a set, would be a finding, and the appendix
contains none. The paper's gate exits zero. And the compiled document carries no
unresolved reference or citation: two independent mechanisms enforce that, since
the third check refuses a document whose cross-references, citation keys or
section labels do not close, and the log scan refuses an output on which the
engine reported one.

One thing is deliberately not automated, and it is the limit of the audit. The
correspondence between a sentence of this paper and a Lean statement is checked
at the level of declaration names and of the frozen digests, never at the level
of the statement's meaning: the name check confirms that a quoted name exists in
the delivered development, and the digest confirms that the statement behind a
frozen name has not changed since it was accepted, but no script compares the
prose of a theorem environment with the type of a declaration, so a sentence
that misdescribes the statement it cites would pass every check we run. That
residual gap is why \cref{sec:lean-boundary} exists as a list of facts rather
than as a claim of completeness.

One episode belongs in this section because it bears on how much a reproducer should
trust the apparatus. An audit of this paper, carried out after it was first
assembled, found that one module of the package was missing the header of a
declaration --- the body survived, the signature did not --- so the sources as
delivered could not have been elaborated, and every claim resting on a build would
have been unreproducible. The header is restored, and the sources have since been
migrated to the Mathlib revision the toolchain pins; the whole package now
elaborates, on that revision, with no error. The transcript in Appendix~B was
\emph{regenerated} afterwards, by \lmod{audit/make_transcript.py}, from the Lean
run on these sources: it therefore records a command that was executed on the
delivered tree, not a transcription of an earlier session. The script also writes
the digest of each declaring module into the transcript header, so a change to any
of them invalidates the record visibly rather than silently.

What that does not buy is stated in the same breath. A footprint is evidence about
the proof term of one declaration; it is not evidence that the declaration says
what the corresponding theorem of Part~I says, and it is not evidence about the
mathematics of any hypothesis in the statement. Those two gaps are the subject of
\cref{sec:lean-boundary} and of the paragraph above, and neither is closed by a
green build.
